\section{Preliminaries}\label{sec:prelims}

We now state formally the main definitions of this paper, and prove some basic equivalences. 

       \paragraph{Notation, graphs.}
       For two distinct elements $a,b$ of a set $V$,
by $ab$ we denote the unordered pair $\set{a,b}$, and call $ab$ a \emph{pair} for simplicity
-- the distinction between ordered pairs and unordered pairs will follow from the context.
% For a set $V$, denote ${V\choose 2}:=\setof{\set{a,b}}{a,b\in V,a\neq b}$.
For $A,B\subset V$, denote $$AB\coloneqq \setof{ab}{a\in A,b\in B,a\neq b}$$ and ${A\choose 2}\coloneqq AA$.
A \emph{graph} is a pair $G=(V,E)$ with $E\subset {V\choose 2}$.
We call the elements of $V$ \emph{vertices}, and elements of $E$ \emph{edges}, and denote $V(G)\coloneqq V$ and $E(G)\coloneqq E$. 
%Denote $\bar E:={V\choose 2}-E$.
%\jan{we use \(\bar E\) only twice. maybe dont define it.}

\subsection{Merge-width and merge sequences}\label{sec:merge-sequence}

In the gray box of \Cref{sec:intro} we gave a definition of merge-width in terms of \emph{construction sequences}
(see also \Cref{sec:computing} for a more formal definition, in the more general setting of binary structures).
We now give an alternative definition in terms of so-called \emph{merge sequences}, and prove in \Cref{lem:fmw}
that these definitions are indeed equivalent.
These two definitions complement each other: 
While construction sequences are useful for showing that a graph class with bounded merge-width has a certain property,
merge sequences are useful for the opposing goal of showing that a given graph class has bounded merge-width.

% Fix a set $V$ of vertices.
% A \emph{merge-sequence} is a sequence of steps, maintaining a partition $\cal P$~of~$V$ and a partition of $V\choose 2$ into three sets: \emph{edges}~$E$, \emph{non-edges} $N$, and \emph{unresolved} pairs $U$. Initially, $\cal P$ partitions $V$ into singletons, and all pairs in $V\choose 2$ are unresolved.
% In each step, we can either \emph{merge} two parts $A,B$ of $\cal P$, replacing them by the single part $A\cup B$, 
% or \emph{resolve} a pair of parts $A,B\in\cal P$ (possibly $A=B$), by proclaiming each of the yet unresolved vertex pairs $\set{a,b}\in U$ with $a\in A$ and $b\in B$ either an edge, or a non-edge -- all in \emph{the same way}. 
% We require that in the end, every pair of vertices in $V$ either forms an edge or a non-edge, thus defining a graph.

% \sz{directly define merge sequences}
% Fix a vertex set $V$.
% A \emph{construction sequence} is a sequence of steps, maintaining a partition $\cal P$~of~$V$ and a partition of  ${V\choose 2}$ into three sets:  
% \emph{edges} $E$, \emph{non-edges} $N$,
% and \emph{unresolved} pairs $U$.
%  Initially, $\cal P$ partitions $V$ into singletons, and every pair in $V\choose 2$ is unresolved.
% In each step, one of three operations is performed:
% \begin{itemize}
%   \item  \emph{merge} two parts of $\cal P$, replacing the two parts by their union $A\cup B$
%   (setting $\cal P:=\cal P-\set{A,B}\cup \set{A\cup B}$), 
%   \item \emph{resolve positively} a pair of parts $A,B\in\cal P$ 
%   (possibly $A=B$), declaring 
%  all the unresolved pairs in $U\cap AB$
%  as edges (setting $E:=E\cup (U\cap AB)$ and $U:=U-AB$),
%  \item \emph{resolve negatively} a pair of parts -- as above, but declaring the unresolved pairs in $U\cap AB$ as non-edges (setting $N:=N\cup (U\cap AB)$ and $U:=U-AB$).
% \end{itemize}
% In the end, we require that every pair in ${V\choose 2}$ is an edge or a non-edge, and that $\cal P$ has one part. This way, we have constructed a graph $G=(V,E)$.

% Fix $r\in\N\cup\set{\infty}$.
% For a set $R\subset {V\choose 2}$ and partition $\cal P$ of $V$,
% the \emph{radius-$r$ width} of $(\cal P,R)$ is the maximum, over all vertices $v\in V$,
% of the number of parts $A\in \cal P$ 
% such that there is a path of length at most $r$ in the graph $G=(V,R)$ from $v$ to some $a\in A$.

% The \emph{radius-$r$ width}
% of a construction sequence is the maximum,
% over all steps in the sequence,
% of the radius-$r$ width of $(\cal P,R)$, where $\cal P$ is the current partition and $R=E\cup N$ 
% is the set of currently resolved pairs (both edges and non-edges). 
% The \emph{radius-$r$ merge-width} of a graph~$G$, denoted $\mw_r(G)$,
% is the least radius-$r$ width of a construction sequence of $G$.
% It is crucial to note that if $\mw_r(G)\le k$, then
%  at any time in the considered merge sequence, and for every part $P\in\cal P$,
% even though we require each vertex $v\in P$ to be adjacent in $(V,E\cup N)$ to at most $k$ of parts of $\cal P$,
% we allow for different vertices in $P$ 
% to be adjacent to different parts.
% If we additionally demanded that throughout the merge sequence, vertices in any given part $P\in\cal P$ are adjacent in $(V,E\cup N)$ to the same parts in $\cal P$, then the resulting parameter for radius $1$
%  would differ from twin-width, $\tww(G)$, by at most one,
% and the parameters for radius $r$ would be upper bounded by $1+\tww(G)+\ldots+\tww(G)^r$.
%  Finally, a graph class $\CC$ has \emph{bounded merge-width}
% if ${\mw_r(\CC)<\infty}$ holds\footnote{For a graph parameter $f$ and graph class $\CC$, by $f(\CC)$ we denote $\sup_{G\in\CC}f(G)$.} for every (finite) $r\in\N$.

\medskip
% The following, alternative perspective, is easier to work with in many contexts.


% , and also not to specify if the performed reveals are positive and which are negative, as this information can be inferred from the target graph $G$.
% This leads to the following, alternative definition of merge-width.

% \begin{definition}
    Fix $r\in\N\cup\set{\infty}$.
        For a set $R\subset {V\choose 2}$ and partition $\cal P$ of $V$,
        we say a part \(A \in \cal P\) is \emph{\(r\)-reachable} from a vertex \(v\) if there is a path of length at most $r$ in the graph $G=(V,R)$ from $v$ to some vertex in~$A$.
        The \emph{radius-$r$ width} of $(\cal P,R)$  is 
        \[
           \max_{v \in V}\ \Big|\bigl\{ A \in \cal P \mid  A \text{ is \(r\)-reachable from } v \bigr\}\Big|.
        \]
        % \end{definition}
        
       


\begin{definition}\label{def:ms}
A \emph{merge sequence} of a graph $G=(V,E)$ 
is a sequence
$$(\cal P_1,R_1),\ldots,(\cal P_m,R_m),$$
  such that:
  \begin{enumerate}
    \item $\cal P_1,\ldots,\cal P_m$ is a sequence of partitions of $V$, 
    where $\cal P_1$ is the partition into singletons and $\cal P_m$ has one part,
    and $\cal P_{t}$ is coarser than (or equal to) $\cal P_{t-1}$ for $t=2,\ldots,m$,
    % \item $\Delta_1,\ldots,\Delta_m$ are 
    % default graphs on $\cal P_1,\ldots,\cal P_m$, respectively,
    \item $R_1\subset\ldots\subset R_m\subset {V\choose 2}$; the pairs in $R_t$ are said to be \emph{resolved} at time $t$, and
    \item for all $t\in [m]$ and $A,B\in \cal P_t$,
     either $AB-R_t\subset E$, or $AB-R_t\subset {V\choose 2} - E$.
  \end{enumerate}
For $r\in\N\cup\set{\infty}$, the \emph{radius-$r$ width} of the merge sequence is the maximum, over $t=2,\ldots,m$,
of the radius-$r$ width of $(\cal P_{t-1},R_t)$.
\end{definition}

Note the offset in the indices of $\cal P_{t-1}$ and $R_t$.
Conceptually, this offset allows us to not only do a single merge, but an unbounded number of merges when going from partition \(\cal P_{t-1}\) to \(\cal P_t\) (see backwards direction of \Cref{lem:fmw}).

\begin{definition}
  The \emph{radius-$r$ merge-width} of a graph $G$, denoted $\mw_r(G)$, is the least radius-$r$ width of a merge sequence of $G$.
Finally, a graph class $\CC$ has \emph{bounded merge-width}
  if ${\mw_r(\CC)<\infty}$ holds\footnote{For a graph parameter $f$ and graph class $\CC$, by $f(\CC)$ we denote $\sup_{G\in\CC}f(G)$.} for every (finite) $r\in\N$.
\end{definition}

\begin{lemma}\label{lem:fmw}
    For every construction sequence of a graph $G$ there is a merge sequence of $G$
  of the same radius-$r$ width, for all $r\in\N\cup\set\infty$,
  and vice versa.
  % For every graph $G$ and $r\in\N\cup\set{\infty}$,
  %    $\mw_r(G)$ is equal to the least radius-$r$ width of a merge sequence for $G$.
  \end{lemma}
  \begin{proof}
    % Fix a  graph $G$. We show that 
    % . This yields the conclusion.
  
  First, fix a construction sequence of $G$, consisting of $m$ steps. 
  % Without loss of generality, we may assume that the merge sequence ends with the partition $\cal P=\set{V}$, as we can always merge parts until only one part remains, without changing the radius-$r$ width of the sequence, for all $r\in\N\cup\set{\infty}$.
    Let $\cal P_t$ be the partition of $V$,
    and let $R_t$ be the set of edges and non-edges 
    at the beginning of step $t$ of the sequence.
    It is clear that 
    $(\cal P_1,R_1),\ldots,(\cal P_m,R_m)$ is a merge sequence. 
    In particular, the third item of \Cref{def:ms} follows from the second item of \Cref{remark:maindef}.
    
    Fix $r\in\N\cup\set{\infty}$, and let $k$ 
    be the radius-$r$ merge-width of the considered construction sequence.
    Then, for each $t\in[1,m]$, 
    the radius-$r$ width of $(\cal P_t,R_t)$ is at most $k$.
    Since for $t>1$, we either have $\cal P_{t-1}=\cal P_t$ or $R_{t-1}=R_t$, it follows that 
    the radius-$r$ width of $(\cal P_{t-1},R_t)$ is at most $k$ as well.
    This concludes one direction.
  
  
  \medskip
  
  In the other direction,  
  let $(\cal P_1,R_1),\ldots,(\cal P_m,R_m)$ be a merge sequence of $G$.
  We define a construction sequence which consists of $m$ stages, as follows.
  In the beginning of stage $t\in[m-1]$,
  the current partition $\cal P$ is the partition $\cal P_t$, and the set $R$ of resolved pairs is contained in $R_t$.
  In stage $t$, first resolve (in any order) all pairs 
  $A,B\in\cal P_{t}$ such that $AB\subset R_{t+1}$.
  The third item of \Cref{def:ms} guarantees that this is always possible by positive or negative resolves.

  Next, repeatedly merge (in any order) any two parts $A,B$ of the current partition, as long as $A\cup B$ is contained in some part of $\cal P_{t+1}$.
  Once this is done, the current partition is equal to $\cal P_{t+1}$, and we proceed to the next stage.
  
  During all steps in stage $t$, the set $R$ of resolved pairs is contained in $R_{t+1}$, and $\cal P$ is a coarsening of $\cal P_t$. In particular, for every $r\in\N\cup\set{\infty}$,
  \begin{multline*}
\text{radius-$r$ width of $(\cal P,R)$}\ \le \ 
\text{radius-$r$ width of $(\cal P_t,R_{t+1})$} \\
\le \ \text{radius-$r$ width of the given merge sequence}.
  \end{multline*}
 
  After the last stage, the current partition $\cal P$ has only the single part $V$. 
  Finally, we resolve $V$ with itself. This completes the  construction sequence.
  \end{proof}

%\[
%   -----
%\]
%\jan{we need normalized construction sequences, not merge sequences.}
%
%
%\medskip
%  Say that a merge sequence for a graph $G$ is \emph{normalized} if in the corresponding sequence $\cal P_1,\ldots,\cal P_m$, for all $t\in [m-1]$, the partition
%  $\cal P_{t+1}$ is obtained from $\cal P_t$ by merging exactly two parts into one. In particular, $m=|V(G)|$.
%
%
%
%  \begin{lemma}\label{lem:normalize}
%    Every merge sequence can be converted to a normalized merge sequence without increasing the radius-$r$ width, for all $r\in\N\cup\set{\infty}$.
%  \end{lemma}
%  \begin{proof}
%    Observe that we can always convert a merge sequence to one in which all partitions $\cal P_1,\ldots,\cal P_m$ are pairwise distinct, since if $\cal P_{t-1}=\cal P_{t}$,
%  then omitting the pair $(\cal P_t,R_t)$ from the merge sequence yields a valid merge sequence of the same radius-$r$ width, for all $r\in\N\cup\set{\infty}$.
%
%  Furthermore, we may assume that for any two consecutive partitions $\cal P_t,\cal P_{t+1}$ of the merge sequence, there is no partition $\cal P$ which is strictly coarser than $\cal P_t$ and strictly finer than $\cal P_{t+1}$. Indeed, otherwise we may insert the pair $(\cal P,R_{t+1})$ between the pairs $(\cal P_t,R_t)$ and $(\cal P_{t+1},R_{t+1})$ of the merge sequence, without increasing its radius-$r$ width, for all $r\in\N\cup\set{\infty}$.
%  \end{proof}

%  \begin{remark}\label{rem:equiv}
%  The merge sequences defined in the introduction (the gray box of \Cref{sec:intro}) correspond to a special form of merge sequences considered in \Cref{def:ms}, namely, merge sequences
%  $(\cal P_1,R_1),\ldots,(\cal P_m,R_m)$ such that
%  for all $t\in[m-1]$,
% either $\cal P_{t+1}=\cal P_t$, or $\cal P_{t+1}$ is obtained from $\cal P_t$ by merging two parts into one, and $R_{t+1}=R_t$.
% Call such merge sequences \emph{step-by-step}.
% For a step-by-step merge sequence, the maximum radius-$r$ width of the pairs $(\cal P_t,R_{t+1})$, for $t\in[m-1]$,
% is equal to the maximum radius-$r$ width of the pairs
% $(\cal P_t,R_t)$, for $t\in[m]$, due to the fact that for all $t\in[m-1]$, either $\cal P_t=\cal P_{t+1}$, or $R_t=R_{t+1}$ and $\cal P_t$ is a refinement of $\cal P_{t+1}$.
% Therefore, the radius-$r$ width of such sequences coincides with the definition given in the introduction.
%
% Furthermore, we can easily convert a normalized merge sequence $(\cal P_1,R_1),\ldots,(\cal P_m,R_m)$ into a step-by-step merge sequence, namely
%$(\cal P_1,R_1),(\cal P_1,R_2),(\cal P_2,R_2),(\cal P_2,R_3),\ldots,(\cal P_m,R_m)$, while preserving the radius-$r$ width of the sequence.
%This, together with \Cref{lem:normalize}, proves the equivalence of the two definitions of merge-width.
%  \end{remark}

  

%  \subsection{Construction sequences}\label{sec:construction-sequences}
%Recall the definition of \emph{construction sequences}, given in the gray box of \Cref{sec:overview}.
%   The following perspective on merge-width is useful in the algorithmic context, as it allows to view merge sequences as 
% a process constructing the target graph.

% \begin{definition}
% Fix a vertex set $V$.
% A \emph{construction sequence} is a sequence of steps, maintaining a partition $\cal P$~of~$V$ and a partition of  ${V\choose 2}$ into three sets:  
% \emph{edges} $E$, \emph{non-edges} $N$,
% and \emph{unresolved} pairs~$U$.
%  Initially, $\cal P$ partitions $V$ into singletons, and every pair in $V\choose 2$ is unresolved.
% In each step, one of three operations is performed:
% \begin{itemize}
%   \item  \emph{merge} two parts $A,B$ of $\cal P$, replacing the two parts by their union $A\cup B$
%   (setting $\cal P:=\cal P-\set{A,B}\cup \set{A\cup B}$), 
%   \item \emph{resolve positively} a pair of parts $A,B\in\cal P$ 
%   (possibly $A=B$), declaring 
%  all the unresolved pairs in $U\cap AB$
%  as edges (setting $E:=E\cup (U\cap AB)$ and $U:=U-AB$),
%  \item \emph{resolve negatively} a pair of parts -- as above, but declaring the unresolved pairs in $U\cap AB$ as non-edges (setting $N:=N\cup (U\cap AB)$ and $U:=U-AB$).
% \end{itemize}
% In the end, we require that every pair in ${V\choose 2}$ is an edge or a non-edge, and that $\cal P$ has one part. This way, we have constructed a graph $G=(V,E)$.

% The \emph{radius-$r$ width}
% of a construction sequence is the maximum,
% over all steps in the sequence,
% of the radius-$r$ width of $(\cal P,R)$, where $\cal P$ is the current partition and $R=E\cup N$ 
% is the set of currently resolved pairs (both edges and non-edges). 
% \end{definition}


  
\subsection{Construction sequences have linear length}\label{sec:representing}

When designing merge-width based algorithms,
we always assume graphs to be represented by construction sequences, rather than merge sequences,
as the former impose more structure.
A~priori such a construction sequence can be arbitrarily long, as for example, the same pair of parts can be resolved indefinitely.
In this subsection we make an innocuous assumption that guarantees a linear bound on their length.


Namely, call a resolve operation in a construction sequence 
\emph{effective} if it creates a non-zero number of new edges or non-edges. 
That is, an \emph{ineffective} resolve operation 
is a resolve operation performed between a pair $A,B$ of parts 
such that all pairs in $AB$ were already resolved before the operation.
Call a construction sequence \emph{effective} if 
every resolve operation in it is effective.
We assume without loss of generality and without explicitly mentioning it that all construction sequences are effective.
Note that the following lemma also applies if the radius-\(r\) width is bounded, as this implies a bound on the radius-\(1\) width.

\begin{lemma}\label{lem:effective-construction-seq}
  Every effective construction sequence of a graph $G$ of radius-$1$ width $w$ has length at most $(2w+1)|V(G)|$.
\end{lemma}
\begin{proof}
  We rearrange the construction sequence by repeatedly postponing resolve operations.
  Namely, consider a subsequence of consecutive operations in a construction sequence,
  consisting of a sequence of resolve operations, and finishing with the merge
  of two parts $A,B\in\cal P$. We may postpone all resolves which occur in the subsequence, and involve 
  neither $A$ or $B$, until after the merge of $A$ and $B$ is done.
  This result in a valid construction sequence,
  of the same radius-$1$ width and the same length.
  As all resolve operations are effective,
  at most $2w$ resolve operations between pairs $P,Q$ remain in the subsequence: at most $w$ pairs involving $A$, and at most $w$ pairs involving $B$, because each vertex in $A$ (or in $B$) $1$-reaches all parts resolved with $A$ (resp. $B$). 
  In the final construction sequence, at most $2w$ resolve operations are performed per merge,
  followed potentially by a single last resolve on the final part.
  As there can be at most \(|V(G)|-1\) merges, we have at most $(2w+1)|V(G)|$ steps altogether.
  Since the original construction sequence has the same length as the rearanged construction sequence, the result follows.
\end{proof}


% \begin{szin}
%     Jan: I updated the following overview.

%     We currently have a  plethora of variants of merge-width:
% \begin{description}
%   \item[construction sequences] Used in the intro and the algorithm.
%    \item[construction sequences for binary structures] Actually used in the proofs.
%   \item[merge sequences] Used for bounded expansion and closure under transductions.
  
%   \item[restrained flip sequences]
%   Used to bound flip-width. Also, to explain the connection with the sparse parameters (wcol, tw), when flips are isolations.
% \end{description}
% Another possible definition is:
% \begin{description}
% \item[merge sequences with default graphs]In each step, we merge two parts \emph{and} define a new default graph on the parts. Essentially, this is a step-by-step merge sequence 
% with radius-$r$ width measured only immediately after each merge. This gives a $\pm 1$ discrepancy with the other definitions, which however slightly  better the twin-width, degeneracy, and clique-width parameters.
% \end{description}


% We should redisucss which definitions we want to include, and also consider including all the definitions (and their equivalence) in the preliminaries section, for a better organization (especially the restrained flip sequences with their obscure discussion).
% \end{szin}


%   \begin{remark}
%     We remark on the equivalence of the two definitons of merge-width -- 
%     If we take a merge sequence, convert it 
%     into a construction sequence 
%     and then back into a merge sequence using the construction in \Cref{lem:fmw}, 
%     we obtain a merge sequence $(\cal P_1,R_1),\ldots,(\cal P_m,R_m)$ such that 
%     for all $t\in[m-1]$,
%     we have that either $R_t=R_{t+1}$ or $\cal P_{t}=\cal P_{t+1}$, and moreover, either $\cal P_{t+1}=\cal P_t$ or $\cal P_{t+1}$ is obtained from $\cal P_t$ by merging two parts into one.
%     This is therefore a merge sequence of the form as considered in the introduction. Moreover, from 
%     the fact that $R_t=R_{t+1}$ or $\cal P_t=\cal P_{t+1}$
%     it follows that the maximum width of $(\cal P_t,R_t)$, for $t\in[m]$,
%     is equal to the maximum width of $(\cal P_{t},R_{t+1})$, for $t\in[m-1]$,
%     so the resulting merge sequence has the same radius-$r$ width as the original one.
%   \end{remark}

  
